\Needspace{7\baselineskip}\subsection{正式入口适配}\label{code:entry}
\noindent 原文件：\path{validation/backend.py}，第 39--48 行。以下逐字摘录，行号保留原编号。
\VerbatimInput[firstnumber=39]{snippets/entry.txt}

\Needspace{7\baselineskip}\subsection{共同几何校验}\label{code:geometry}
\noindent 原文件：\path{shared/geometry.py}，第 39--62 行。以下逐字摘录，行号保留原编号。
\VerbatimInput[firstnumber=39]{snippets/geometry.txt}

\Needspace{7\baselineskip}\subsection{AG 活动状态重建与配对面积演化}\label{code:state}
\noindent 原文件：\path{AG/beam_dynamics_6d.py}，第 1024--1051 行。以下逐字摘录，行号保留原编号。
\VerbatimInput[firstnumber=1024]{snippets/state.txt}

\Needspace{7\baselineskip}\subsection{AG 三方向面积变化率}\label{code:g}
\noindent 原文件：\path{AG/beam_dynamics_6d.py}，第 356--369 行。以下逐字摘录，行号保留原编号。
\VerbatimInput[firstnumber=356]{snippets/g.txt}

\Needspace{7\baselineskip}\subsection{AG 高斯 SC 纵向力及其上下文}\label{code:force}
\noindent 原文件：\path{AG/beam_dynamics_6d.py}，第 318--340 行。以下逐字摘录，行号保留原编号。
\VerbatimInput[firstnumber=318]{snippets/force.txt}

\Needspace{7\baselineskip}\subsection{直接积分及域外保护}\label{code:alpha}
\noindent 原文件：\path{AG/alpha_coefficients.py}，第 51--100 行。以下逐字摘录，行号保留原编号。
\VerbatimInput[firstnumber=51]{snippets/alpha.txt}

\Needspace{7\baselineskip}\subsection{共同 RF 系数与正弦调制}\label{code:rfshared}
\noindent 原文件：\path{shared/rf_model.py}，第 11--27 行。以下逐字摘录，行号保留原编号。
\VerbatimInput[firstnumber=11]{snippets/rfshared.txt}

\Needspace{7\baselineskip}\subsection{AG 协方差薄透镜}\label{code:rfag}
\noindent 原文件：\path{AG/beam_dynamics_6d.py}，第 783--814 行。以下逐字摘录，行号保留原编号。
\VerbatimInput[firstnumber=783]{snippets/rfag.txt}

\Needspace{7\baselineskip}\subsection{OCELOT 实际传输方法守门}\label{code:method}
\noindent 原文件：\path{GPT模拟/ued_beamline_v2.py}，第 195--209 行。以下逐字摘录，行号保留原编号。
\VerbatimInput[firstnumber=195]{snippets/method.txt}

\Needspace{7\baselineskip}\subsection{OCELOT 原生 SC 调度}\label{code:schedule}
\noindent 原文件：\path{GPT模拟/ued_beamline_v2.py}，第 518--536 行。以下逐字摘录，行号保留原编号。
\VerbatimInput[firstnumber=518]{snippets/schedule.txt}

\Needspace{7\baselineskip}\subsection{共同中心矩与纵向输出}\label{code:diagnostic}
\noindent 原文件：\path{shared/diagnostics.py}，第 22--42 行。以下逐字摘录，行号保留原编号。
\VerbatimInput[firstnumber=22]{snippets/diagnostic.txt}

\Needspace{7\baselineskip}\subsection{GUI 调用正式模型}\label{code:worker}
\noindent 原文件：\path{gui/worker.py}，第 36--70 行。以下逐字摘录，行号保留原编号。
\VerbatimInput[firstnumber=36]{snippets/worker.txt}

\Needspace{7\baselineskip}\subsection{Windows 父进程存活检查（尚待实机验收）}\label{code:platform}
\noindent 原文件：\path{gui/platform_compat.py}，第 170--218 行。以下逐字摘录，行号保留原编号。
\VerbatimInput[firstnumber=170]{snippets/platform.txt}
